\documentclass[11pt,a4paper]{article}

\usepackage[utf8]{inputenc}
\usepackage[T1]{fontenc}
\usepackage[portuguese]{babel}
\usepackage{helvet}
\renewcommand{\familydefault}{\sfdefault}
\usepackage[top=1.8cm, bottom=1.8cm, left=2.5cm, right=2.5cm]{geometry}
\usepackage{setspace}
\usepackage{booktabs}
\usepackage{array}
\usepackage{microtype}
\usepackage{parskip}
\usepackage{titlesec}
\usepackage{bm}

\setstretch{1.15}

\titleformat{\section}{\normalfont\bfseries\normalsize}{\thesection.}{0.5em}{}
\titlespacing{\section}{0pt}{5pt}{2pt}

\setlength{\parskip}{3pt}
\setlength{\parindent}{0pt}

\begin{document}

\begin{center}
    {\large\bfseries Relatório Analítico: Algoritmo PageRank}\\[2pt]
    {\small Nome: Deivid Mark Brayan Gomes Andrade \quad|\quad Matrícula: 2023093222}
\end{center}

\vspace{3pt}
\hrule
\vspace{5pt}

\section{Comparação de Métodos: Amostragem vs.\ Método da Potência}

Ambos os métodos produziram resultados convergentes, conforme esperado pela teoria. A tabela abaixo sintetiza os valores obtidos:

\begin{center}
\small
\setstretch{1.0}
\begin{tabular}{lcc}
\toprule
\textbf{Página} & \textbf{Amostragem} & \textbf{Iteração} \\
\midrule
portal\_gamer.html  & 0.2246 & 0.2202 \\
guia\_rpg.html      & 0.2040 & 0.1975 \\
forum\_gamer.html   & 0.1921 & 0.1975 \\
guia\_mmorpg.html   & 0.1137 & 0.1145 \\
discord\_gamer.html & 0.1124 & 0.1145 \\
review\_gta.html    & 0.0797 & 0.0778 \\
review\_zelda.html  & 0.0735 & 0.0778 \\
\bottomrule
\end{tabular}
\end{center}

As divergências (máx.\ $\approx 0{,}005$) são atribuídas à variância residual $O(1/\sqrt{n})$ da amostragem por Random Walk com $n=10{.}000$. À luz da Lei dos Grandes Números, as simetrias exatas do método iterativo --- como \texttt{guia\_rpg} e \texttt{forum\_gamer} em 0.1975 --- não se manifestam perfeitamente na amostragem devido a esse ruído estocástico. Em essência, ambos os métodos convergem para o mesmo vetor estacionário $\mathbf{r}$ da cadeia de Markov definida pela matriz de transição, validando a corretude das implementações.

\section{Impacto do Fator de Amortecimento}

O fator $d$ controla a mistura entre a navegação por links e a teleportação aleatória. Na equação iterativa $\mathbf{r} = d \cdot M\mathbf{r} + \frac{1-d}{N}\mathbf{1}$, ele determina o raio espectral da matriz modificada: quanto maior o valor de $d$, mais lenta é a convergência, pois o autovalor dominante se aproxima de 1.

\textbf{Se $d \to 0$:} a teleportação domina e o vetor $\mathbf{r}$ converge para $1/N$ uniformemente em todas as páginas — a topologia da rede torna-se irrelevante. \textbf{Se $d \to 1$:} a teleportação desaparece e o modelo fica suscetível a armadilhas topológicas como ciclos isolados ou nós sorvedouros (\textit{dangling nodes}), gerando instabilidade ou não convergência. O valor canônico $d = 0{,}85$ preserva a influência topológica, garantindo a irredutibilidade e a ergodicidade da cadeia.

\section{Análise do Corpus: Mini Web Gamer}

A página \texttt{portal\_gamer.html} obteve o maior PageRank (0.22) por atuar como o principal \textit{hub} do sistema, concentrando múltiplos links de entrada vindos de páginas de alta relevância. Sob a ótica da álgebra linear, o PageRank mapeia o autovetor dominante de $M$: nós que recebem fluxo de páginas já bem ranqueadas herdam e acumulam essa importância por propagação espectral.

Os ciclos isolados (\texttt{guia\_rpg $\leftrightarrow$ guia\_mmorpg} e \texttt{forum $\leftrightarrow$ discord}) retêm e redistribuem rank predominantemente entre si, o que justifica seus valores moderados e simétricos. Por fim, a página \texttt{review\_gta.html}, configurando-se como um \textit{dead end}, absorve rank do sistema mas, por não possuir links de saída, é incapaz de propagar sua influência diretamente, dependendo exclusivamente da teleportação aleatória para redistribuir sua massa estocástica.

\end{document}